\section{The finite integral comparison error on growing hosts}
\label{sec:finite-gluing-budget}

The integral Fourier comparison in ALP\_MACHINE, AM-H06,
leaves finite-coefficient cohomology terms.  These terms admit
a uniform elementary bound.  Their control does not supply
the main infinite-coefficient strict Selmer lengths.

Write $n_F=[F:\mathbb Q]$, $D_F=|\operatorname{disc}F|$,
and $s_F=|S_{\mathrm{fin}}|$.  Throughout this section $F$
is totally imaginary and $S$ contains the archimedean places
and all places over the coefficient prime $\ell$.
All logarithms without a subscript are natural.

\begin{theorem}[AC-K19: finite gluing errors fit the growing-host budget]
\label{AC-K19}
Suppose first that $\mu_\ell\subset F$.  Define
\[
 \begin{split}
 d_i(F,S)&=\dim_{\mathbb F_\ell}H^i(G_{F,S},\mathbb F_\ell)
                  \quad(i=1,2),\\
 d_{\mathrm{loc}}(F,S)&=
 \sum_{v\mid\ell}\dim_{\mathbb F_\ell}H^1(F_v,\mathbb F_\ell)
 +\sum_{v\in S_{\mathrm{fin}},\ v\nmid\ell}
                \dim_{\mathbb F_\ell}H^1(I_v,\mathbb F_\ell).
 \end{split}
\]
Then these dimensions are finite and
\[
 d_1+d_2+d_{\mathrm{loc}}
 \leq B_\ell(F,S):=
 2\log_\ell D_F+(3+2\log_\ell8)n_F+5s_F.       \tag{FG.1}
\]

More generally, let $M$ be a finite $G_{F,S}$-module,
killed by $\ell^a$ and generated by at most $r$ elements
as an abelian group.  Let $L/F$ be a finite Galois extension
unramified outside $S$, trivializing $M$ and containing
$\mu_{\ell^a}$.  Put $D=[L:F]$, and let $S_L$ be the
places above $S$.  The quantity
\[
 \begin{split}
 \mathcal C_{F,S}(M)={}&\log|H^1(G_{F,S},M)|
      +\log|H^2(G_{F,S},M)|\\
 &+\sum_{v\mid\ell}\log|H^1(F_v,M)|
      +\sum_{v\in S_{\mathrm{fin}},\ v\nmid\ell}
                         \log|H^1(I_v,M)|
 \end{split}
\]
is finite and satisfies
\[
 \begin{split}
 \mathcal C_{F,S}(M)\leq{}&2r\log D_L+2rn_L\log8\\
 &+ar\log\ell\,[2(D+n_L+s_L)+n_F+3s_F].
 \end{split}
 \tag{FG.2}
\]
In particular, a fixed finite Artin coefficient, with an
arbitrary integral Tate twist, has
$\mathcal C=O(\log D_F+n_F+s_F)$ when trivialized after
compositum with one fixed number field; the constant is
independent of the twist.

For the growing signed hosts of AM-H17, adjoin the cubic
coefficient field and write
\[
 F=\mathbb Q(\zeta_{2^{h+2}},\zeta_7,
               \sqrt{-q},\sqrt{-d},\omega),
 \qquad F^{\mathrm c}=F(\sqrt[3]{269}).          \tag{FG.3}
\]
Here $q,d>7$ are the distinct primes in the source's
$\operatorname{WinPair}(\varepsilon,Y_0)$, $x\geq Y_0\geq\exp(1)$,
and $h$ satisfies its $\operatorname{PageN}$ condition
\[
 x^{2+8\varepsilon}\leq2^h\leq2x^{2+8\varepsilon},
 \qquad 0<\varepsilon<1/52.
\]
Take $S_{\mathrm{fin}}$ to be the primes of $F$ over
$3\cdot269$ and $t=\lfloor x^{401}\rfloor$.
For a positive weight $j$, let $a_{F^{\mathrm c}}(j)$
and $b_F(j)$ be the two sides of AM-H06: the untwisted
strict $\mathbb Z_3$-length over $F^{\mathrm c}$, and the
sum of the three character lengths over $F$, respectively.
Then
\[
 \begin{split}
 \sum_{1\leq j<t}|2a_{F^{\mathrm c}}(j)-b_F(j)|\log3
 &\leq3(t-1)B_3(F,S)\log3\\
 &=O_\varepsilon\!\left(
 x^{403+41\varepsilon/5}(\log x)^2\right)\\
 &=o\!\left(e^{x^\varepsilon/2}\right).
 \end{split}
 \tag{FG.4}
\]
Thus this finite integral comparison error consumes less
than any prescribed positive fraction of the source budget
for sufficiently large $x$.  Absorbing it into an exact
budget requires that positive relative slack in the
main arithmetic estimate.  No such main estimate is
asserted by this theorem.
\end{theorem}

\begin{proof}
\proofstep{1. A coarse class-number bound.}
Let $H_F$ be the ordinary ideal class number.
Minkowski's bound supplies an integral ideal in every
class of norm at most
\[
 \frac{n_F!}{n_F^{n_F}}
       \left(\frac4\pi\right)^{r_2(F)}\sqrt{D_F}
       \leq2^{n_F}\sqrt{D_F}=:U_F.
\]
For the source statement, see Milne,
\emph{Algebraic Number Theory}, v3.08 (2020),
Theorem~4.3, p.~70.
The Euler factors give $\zeta_F(2)\leq\zeta(2)^{n_F}<2^{n_F}$:
at each rational prime there are at most $n_F$ factors,
and every residue degree is at least one.
Counting all ideals of norm at most $U_F$ therefore gives
\[
 H_F\leq U_F^2\sum_{\mathfrak a\ne0}(N\mathfrak a)^{-2}
             \leq8^{n_F}D_F.                         \tag{FG.5}
\]
This estimate uses no hypothesis about zeros.

\proofstep{2. Constant finite coefficients.}
Choose an identification $\mathbb F_\ell\simeq\mu_\ell$
over $F$.  Kummer theory gives
\[
 0\longrightarrow\mathcal O_{F,S}^{\times}/
                    (\mathcal O_{F,S}^{\times})^\ell
 \longrightarrow H^1(G_{F,S},\mu_\ell)
 \longrightarrow\operatorname{Cl}(\mathcal O_{F,S})[\ell]
 \longrightarrow0.                                  \tag{FG.6}
\]
Indeed, outside $S$ the residue characteristic differs
from $\ell$, and a Kummer class is unramified exactly
when its valuation is a multiple of $\ell$.
Dividing its divisor outside $S$ by $\ell$ gives the
right-hand map.  A trivial divisor class leaves an
$S$-unit, and any $\ell$-torsion ideal class lifts by
choosing a generator of its $\ell$th power.  This proves
the displayed exact form from field Kummer theory.

The $S$-unit rank is $r_1(F)+r_2(F)-1+s_F$; its torsion
subgroup is cyclic.  Hence its quotient by $\ell$th
powers has dimension at most $n_F+s_F$.
Here the known unit theorem is Milne, loc. cit.,
Theorem~5.11, p.~90.
The $S$-class group is a quotient of the ordinary class
group, so
\[
             d_1\leq n_F+s_F+\log_\ell H_F.           \tag{FG.7}
\]

The last part of the finite Poitou--Tate sequence is
\[
 H^1(G_{F,S},\mu_\ell)^\vee
 \longrightarrow H^2(G_{F,S},\mathbb F_\ell)
 \longrightarrow\bigoplus_{v\in S}H^2(F_v,\mathbb F_\ell).
\]
Finite local duality bounds the dimension of each finite-place
term on the right by one; complex places contribute zero.
Consequently $d_2\leq d_1+s_F$.
The exact sources are Sharifi, \emph{Iwasawa Theory},
Appendix~A, Theorems~A.0.3 and A.0.8; these are the
finite-module statements, with $S$ containing all places
over $\ell$ and infinity.

Local duality and the local Euler characteristic give,
for $v\mid\ell$,
\[
      \dim H^1(F_v,\mathbb F_\ell)=[F_v:\mathbb Q_\ell]+2.
\]
The Euler-characteristic source is Milne,
\emph{Arithmetic Duality Theorems}, second edition (2006),
I, Theorem~2.8, p.~31.
If $v\nmid\ell$, wild inertia is pro-$p$ with $p\ne\ell$
and tame inertia has maximal pro-$\ell$ quotient
$\mathbb Z_\ell(1)$.  Thus
$\dim H^1(I_v,\mathbb F_\ell)=1$.
Summing the local degrees yields
\[
 d_{\mathrm{loc}}\leq n_F+2s_F.
\]
Combining this with (FG.7), $d_2\leq d_1+s_F$, and
(FG.5) proves (FG.1).

\proofstep{3. Arbitrary bounded finite coefficients.}
Set $M^D=\operatorname{Hom}(M,\mu_{\ell^a})$.
It has the same elementary divisors as $M$, is trivial
over $L$, and is the finite Tate dual used by duality.
Inflation--restriction and the elementary bound on
one-cocycles of the finite group $\operatorname{Gal}(L/F)$
give
\[
 |H^1(G_{F,S},M)|
 \leq |M|^D\,|H^1(G_{L,S_L},M)|.
\]
Over $L$, decompose $M$ into at most $r$ cyclic groups
of orders $\ell^{b_i}$, where $b_i\leq a$.
Since $\mu_{\ell^a}\subset L$, the same Kummer exact
sequence as (FG.6), now with exponent $\ell^{b_i}$,
bounds each factor by
$\ell^{b_i(n_L+s_L)}H_L$.  Therefore, for both $M$ and
$M^D$,
\[
 \log|H^1(G_{F,S},M)|
 \leq ar\log\ell\,(D+n_L+s_L)+r\log H_L.           \tag{FG.8}
\]
Finite Poitou--Tate and local duality give
\[
 \log|H^2(G_{F,S},M)|
 \leq\log|H^1(G_{F,S},M^D)|+s_F\log|M|.           \tag{FG.9}
\]
For $v\mid\ell$, the same local Euler characteristic
and duality imply
$|H^1(F_v,M)|\leq|M|^{[F_v:\mathbb Q_\ell]+2}$.
For $v\nmid\ell$, first taking wild-inertia invariants
and then the invariants of the prime-to-$\ell$ part of
tame inertia is exact on finite $\ell$-groups.
The remaining first cohomology is a quotient
$N/(\tau-1)N$ for a submodule $N$ of $M$ and a
generator of $\mathbb Z_\ell$.  It has order at most
$|M|$.  The two local sums in $\mathcal C$ are thus at
most $(n_F+2s_F)\log|M|$.
Substitute (FG.8)--(FG.9), $|M|\leq\ell^{ar}$ and
(FG.5) for $L$.  This is (FG.2).

For the fixed-field corollary, take a fixed finite
Galois field $L_0$ trivializing the Artin action and
containing $\mu_{\ell^a}$.  It trivializes all Tate
twists as well.  Put $L=FL_0$, and include the fixed
ramified primes in $S$.  If $C=[L_0:\mathbb Q]$, then
$D\leq C$, $n_L\leq Cn_F$ and $s_L\leq Cs_F$.
Moreover
\[
             D_L\leq D_F^C D_{L_0}^{n_F}.            \tag{FG.10}
\]
For a direct proof valid even when the fields intersect,
the finite etale algebra $F\otimes_{\mathbb Q}L_0$ is
a product of number fields, one of them $L$.
The tensor order $\mathcal O_F\otimes\mathcal O_{L_0}$
has discriminant $D_F^C D_{L_0}^{n_F}$ in absolute value,
by the determinant of the tensor trace matrix.
The product maximal order has discriminant the product
of the positive absolute field discriminants; this
divides the tensor-order discriminant by the square
of the index.  Each factor is at least one, proving
(FG.10).  Equation (FG.2) now has the stated uniform order.

\proofstep{4. The actual cubic growing hosts.}
The five fields in (FG.3) have pairwise coprime
conductors $2^{h+2},7,q,d,3$.
They are linearly disjoint: their Galois groups occur
in distinct factors of the cyclotomic Chinese-remainder
product, and the quadratic fields use the $q$ and $d$
factors.  Their discriminants are also pairwise coprime.
Milne, \emph{Algebraic Number Theory}, Lemma~6.5,
p.~100, and Remark~6.6(c), p.~101, identify the product
integral basis and its discriminant.  The cyclotomic
discriminant formula in that remark gives
\[
 \begin{split}
 n_F&=96\cdot2^h,\\
 \log D_F/n_F&=(h+1)\log2+\tfrac56\log7
                 +\tfrac12\log q+\tfrac12\log d
                 +\tfrac12\log3.
 \end{split}
 \tag{FG.11}
\]
The quadratic discriminants are $-q$ and $-d$, because
the source congruences give $q,d\equiv3\pmod4$.
There are at most $2n_F$ places in $S_{\mathrm{fin}}$.
The field $F$ is abelian over $\mathbb Q$, so it cannot
contain the nonabelian splitting field
$\mathbb Q(\omega,\sqrt[3]{269})$.
Thus $F^{\mathrm c}/F$ is indeed cyclic of degree three.

The proof of AM-H06 applies over any base containing
$\mu_3$; its local conditions explicitly retain inertia
at $269$.  With $c_3=3$, its equation (H06.3) gives
\[
 -3(d_2+d_{\mathrm{loc}})
 \leq2a_{F^{\mathrm c}}(j)-b_F(j)\leq3d_1.
\]
For positive weights the untwisted group is finite by
AC-K10, and AM-H06 then makes all the displayed lengths
finite.  Equation (FG.1) bounds the absolute difference
by $3B_3(F,S)$, proving the first line of (FG.4).

The exact source window inequalities are
\[
 \log q\leq Y_0^{\varepsilon/5},\qquad
 \log d\leq Y_0^{\varepsilon/5}(\log Y_0)^2.
\]
For $x\geq Y_0\geq e$ these imply the same upper bounds
with $x$ in place of $Y_0$; also $h=O_\varepsilon(\log x)$.
Using $n_F\leq192x^{2+8\varepsilon}$ in (FG.11) yields
\[
 B_3(F,S)=O_\varepsilon\!\left(
 x^{2+8\varepsilon+\varepsilon/5}(\log x)^2\right).
\]
There are fewer than $x^{401}$ weights.
This proves the polynomial estimate in (FG.4).
Finally, for every fixed $A$ and fixed $\varepsilon>0$,
$A\log x+2\log\log x-x^\varepsilon/2\to-\infty$.
The last assertion of (FG.4) and the statement about
arbitrary positive relative slack follow.
\end{proof}

\paragraph{Application to the full finite-character decomposition.}
The two defects of AC-G30-Fourier satisfy (FG.2) with
$(\ell,a,r)=(3,1,24)$ and $(2,3,42)$, respectively.
For the second count, the $21$ nonunit Smith factors
over the unramified quadratic coefficient ring each
have two generators over $\mathbb Z_2$.
The common trivializing field is
$L_0=\mathbb Q(\zeta_{168})$, of degree $48$.
On the base $B_h$ of AC-G30 include the fixed primes
$2,3,7$ in $S$, and also $269$ for the cubic comparison.
The degree and log discriminant of $B_h$ are bounded
by those of $F$, by the tower discriminant formula.
Thus (FG.2), (FG.10) and (FG.11) give the same
polynomial order in (FG.4) for the sum of these two
finite defects over all the weights.  The long exact
global and local coefficient sequences bound the
kernel and cokernel of strict localization exactly as
in AM-H06.  The normalized lengths times $\log\ell$
are at most the logarithms of the finite group orders,
so the bound also applies to those comparison errors.
This controls the full finite Fourier cost without
assuming trivial residual characters on $B_h$; it
does not synchronize the main character windows in
AC-G31.

\paragraph{Source interfaces and remaining gate.}
The exact inherited comparison is
\path{RETURN_ALP_MACHINE_checkpoint_04_20261121T2304Z/parts/hecke/section.tex},
AM-H06, equations (H06.2)--(H06.4); its proof was read
in full for this application.  The window and host
quantifiers are those of
\path{Development/AutomorphicArithmetic/CharacterBoxes/ExistsPageN.lean},
declarations \texttt{WinPair}, \texttt{PageN},
\texttt{VVle}, \texttt{VAt} and \texttt{VWinAt}.
The source's previously unbounded $d_i,d_{\mathrm{loc}}$
are now quantitatively controlled on these hosts.
Neither (FG.1) nor (FG.2) bounds the length of an
infinite-coefficient strict group from its residual
cohomology.  That main depth estimate, its all-prime
sum, and the arithmetic theta tie remain separate.

\paragraph{Algebraic-geometric dictionary and source access.}
The finite module $M$ is a locally constant finite
etale coefficient on $\operatorname{Spec}\mathcal O_F[S^{-1}]$;
the finite covering that trivializes it belongs to the
Galois-category dictionary of SGA~1, Expos\'e~V,
\S\S4 and 7.  The Fourier cokernel is a finite integral
coefficient supported at the bad coefficient primes.
Equations (FG.1)--(FG.4) bound the cost of that coefficient
in the global and local cohomology of the arithmetic
curve.  They do not produce a surface square or a
polarization of the infinite-coefficient object.
No new Grothendieck archive scan is used here; the
archive access qualifications in Section~\ref{sec:groth-source}
remain unchanged.  The primary source URLs for the
number-theoretic inputs are
\url{https://www.jmilne.org/math/CourseNotes/ANT.pdf},
\url{https://www.jmilne.org/math/Books/ADTnot.pdf}, and
\url{https://math.ucla.edu/~sharifi/notes/iwasawa-app-a.html}.
